\documentclass[11pt,a4paper]{article}
\usepackage[margin=1in]{geometry}
\usepackage{amsmath,amssymb,amsthm}
\usepackage{algorithm}
\usepackage{algpseudocode}
\usepackage{booktabs}
\usepackage{hyperref}

\theoremstyle{definition}
\newtheorem{definition}{\textbf{Definition}}[section]
\newtheorem{theorem}{\textbf{Theorem}}[section]
\newtheorem{lemma}[theorem]{\textbf{Lemma}}
\newtheorem{remark}{\textbf{Remark}}[section]

\title{\textbf{DoReMi: Domain Reweighting via Distributionally Robust Optimization and Excess Loss}}
\author{\textbf{Team Scaibu} \\ \small Email: \texttt{jaydeep.9664920749@gmail.com} \\ \small \textbf{Architecture and Core Engine Team}}
\date{\textbf{October 2026}}

\begin{document}
\maketitle

\section{Definitions and Cognitive Mental Models}

\subsection{Formal Mathematical Definition}
\textbf{Definition (DoReMi Domain Reweighting with Excess Loss Optimization):}
Let $\mathcal{D} = \{1, \dots, D\}$ represent a collection of heterogeneous pre-training data domains (e.g., Code, Wikipedia, ArXiv, Web Crawl, Books). A language model trained under domain sampling weights $\boldsymbol{\alpha} = (\alpha_1, \dots, \alpha_D) \in \Delta^D$, where $\Delta^D = \{ \boldsymbol{\alpha} \in \mathbb{R}_+^D : \sum_{d=1}^D \alpha_d = 1 \}$, achieves empirical cross-entropy loss $L_d(\theta)$ on domain $d$. Standard Group Distributionally Robust Optimization (DRO) minimizes worst-case domain loss $\max_d L_d(\theta)$, which suffers from pathological domain starvation when inherently noisy or high-entropy domains (e.g., uncurated web text) monopolize training compute without yielding linguistic utility.

\textbf{DoReMi} (Domain Reweighting with Minimax Optimization) remedies this failure mode by training a small proxy model $\theta_{\text{proxy}}$ against a pre-trained reference baseline model $\theta_{\text{ref}}$, optimizing against \textbf{excess loss} $\Delta_d$:
{\boldmath
\begin{equation}
\Delta_d \triangleq L_d(\theta_{\text{proxy}}) - L_d(\theta_{\text{ref}})
\end{equation}
}
The domain weights are dynamically updated through exponentiated gradient mirror descent over the simplex $\Delta^D$:
{\boldmath
\begin{equation}
\alpha_d^{(t+1)} = \frac{\alpha_d^{(t)} \exp\left( \eta \, \Delta_d^{(t)} \right)}{\sum_{k=1}^D \alpha_k^{(t)} \exp\left( \eta \, \Delta_k^{(t)} \right)}
\end{equation}
}
where $\eta > 0$ denotes the dual learning rate. To prevent domain starvation and ensure non-vanishing exploration, an $\epsilon$-smoothing floor is applied: $\boldsymbol{\alpha}^{(t+1)} \leftarrow (1 - D\epsilon)\boldsymbol{\alpha}^{(t+1)} + \epsilon \mathbf{1}$.

\subsection{Expanded Non-Technical Definition (Intuition for Anyone)}
\begin{enumerate}
    \item \textbf{The Baseline-Calibrated Tutor}: Imagine studying for an exam containing 5 subjects. Calculus is naturally much harder to get a 100\% in than History. If your study schedule only focused on the subject with the lowest raw percentage, you would spend 95\% of your time on Calculus and fail History. DoReMi compares your current grade against an expert student's grade: if you are failing Chemistry by 20 points compared to the expert, but only 5 points in Calculus, DoReMi assigns more study time to Chemistry.
    \item \textbf{Preventing Noise Monopolization}: Some web data is pure noise (garbled HTML strings). A model will always have a high prediction loss on random strings. Standard worst-case optimization gets tricked into studying garbage forever; DoReMi realizes the reference model also struggles with garbage, so excess loss is near zero, and garbage gets downweighted!
    \item \textbf{Cheap Proxy, Expensive Scaling}: Instead of tuning domain ratios on a massive, trillion-parameter model costing millions of dollars, DoReMi tunes the weights using a tiny 200M parameter proxy model over a few hours, then freezes the optimal domain proportions for the final frontier training run.
\end{enumerate}

\subsection{Child-Friendly Mental Model (Physical Metaphor)}
Imagine an orchestra with violinists, drummers, and flutists:
\begin{itemize}
    \item \textbf{Naive Balance}: The conductor tells the drums to play louder simply because drums are loud and hard to tune, drowning out the delicate flutes.
    \item \textbf{DoReMi Conductor}: The conductor compares each musician against a master recording. Whichever section is deviating most from its ideal master harmony receives extra coaching, ensuring every instrument blends into a balanced, rich symphony!
\end{itemize}

\section{Core Mathematical Equations and Definitions}

\subsection{Minimax Excess Loss Formulation}
{\boldmath
\begin{equation}
\min_{\theta \in \Theta} \max_{\boldsymbol{\alpha} \in \Delta^D} \sum_{d=1}^D \alpha_d \left( L_d(\theta) - L_d(\theta_{\text{ref}}) \right) - \frac{1}{\eta} \operatorname{KL}\left(\boldsymbol{\alpha} \,\|\, \boldsymbol{\alpha}_0\right)
\end{equation}
}
\begin{itemize}
    \item \textbf{Formal Definition}: Regularized saddle-point game between model parameters $\theta$ and domain adversarial weighting distribution $\boldsymbol{\alpha}$.
    \item \textbf{What It Means}: The proxy learner attempts to minimize excess loss across all domains while the dual player upweights domains where the proxy lags furthest behind the reference model.
    \item \textbf{Why It Is Important}: Guarantees Pareto-optimal sample allocation without overfitting to domains with irreducible Bayesian entropy.
\end{itemize}

\subsection{Exponentiated Gradient Mirror Descent Update}
{\boldmath
\begin{equation}
\tilde{\alpha}_d^{(t+1)} = \alpha_d^{(t)} \cdot \exp\left( \operatorname{clip}\left(\eta \cdot \Delta_d, -c, c\right) \right), \quad \alpha_d^{(t+1)} = \frac{\tilde{\alpha}_d^{(t+1)}}{\sum_{k=1}^D \tilde{\alpha}_k^{(t+1)}}
\end{equation}
}
\begin{itemize}
    \item \textbf{Formal Definition}: Simplex-constrained Bregman projection with entropy generating function.
    \item \textbf{What It Means}: Computes the multiplicative domain probability update step with numerical overflow guarding.
    \item \textbf{Why It Is Important}: Ensures all weights remain strictly non-negative and properly normalized on the probability simplex $\Delta^D$ with $O(D)$ time complexity.
\end{itemize}

\section{Rigorous Mathematical Analysis Checklist}

\begin{itemize}
    \item \textbf{Notation}: $D$ (number of domains), $\boldsymbol{\alpha} \in \Delta^D$ (domain probability vector), $L_d$ (cross-entropy loss), $\Delta_d$ (excess loss), $\eta$ (dual step size), $\epsilon$ (minimum floor).
    \item \textbf{Epistemic Status}: Proven in empirical foundation model training (Xie et al., 2023) and adopted in standard pre-training pipelines.
    \item \textbf{Dependency Graph}: Proxy \& Reference Losses $\to$ Excess Loss Vector $\to$ Exponentiated Multiplicative Step $\to$ Simplex Normalization $\to$ Domain Weights.
    \item \textbf{Dimensions}: Loss vectors in $\mathbb{R}^D$; weight simplex in $\mathbb{R}^D$ where $\sum_{d=1}^D \alpha_d = 1$.
    \item \textbf{Regimes}: Applicable to large multi-corpus datasets spanning code, science, dialog, web crawl, and structured math.
    \item \textbf{Invariants}: Strict non-negativity ($\alpha_d \ge \epsilon$) and unit simplex conservation ($\sum_{d=1}^D \alpha_d = 1.0$) for all iterations $t$.
    \item \textbf{Isomorphisms}: Isomorphic to Online Mirror Descent with Entropic Regularization in zero-sum matrix games.
\end{itemize}

\section{Theoretical Theorems and Formal Proofs}

\begin{theorem}[Convergence of Exponentiated Gradient Dual Update on Excess Loss]
Let the excess loss vector $\boldsymbol{\Delta}^{(t)} \in \mathbb{R}^D$ be uniformly bounded such that $\|\boldsymbol{\Delta}^{(t)}\|_\infty \le M$ for all $t \in \{1, \dots, T\}$. Under the exponentiated gradient mirror descent update $\alpha_d^{(t+1)} \propto \alpha_d^{(t)} \exp(\eta \Delta_d^{(t)})$ initialized at the uniform distribution $\boldsymbol{\alpha}^{(1)} = \left(\frac{1}{D}, \dots, \frac{1}{D}\right)$ with step size $\eta = \sqrt{\frac{2 \ln D}{M^2 T}}$, the dual regret with respect to any static comparison distribution $\boldsymbol{\alpha}^* \in \Delta^D$ satisfies:
{\boldmath
\begin{equation}
R_T \triangleq \sum_{t=1}^T \langle \boldsymbol{\alpha}^* - \boldsymbol{\alpha}^{(t)}, \, \boldsymbol{\Delta}^{(t)} \rangle \le M \sqrt{2 T \ln D}
\end{equation}
}
\end{theorem}

\begin{proof}
Let $\Phi(\boldsymbol{\alpha}) = \sum_{d=1}^D \alpha_d \ln \alpha_d$ denote the negative Shannon entropy on $\Delta^D$. The Bregman divergence generated by $\Phi$ is the Kullback-Leibler divergence:
\begin{equation}
D_{\Phi}(\boldsymbol{\alpha}^*, \boldsymbol{\alpha}) = \operatorname{KL}(\boldsymbol{\alpha}^* \,\|\, \boldsymbol{\alpha}) = \sum_{d=1}^D \alpha_d^* \ln \frac{\alpha_d^*}{\alpha_d}
\end{equation}
By the standard Mirror Descent potential lemma under entropic mirror map:
\begin{equation}
D_{\Phi}(\boldsymbol{\alpha}^*, \boldsymbol{\alpha}^{(t+1)}) - D_{\Phi}(\boldsymbol{\alpha}^*, \boldsymbol{\alpha}^{(t)}) \le -\eta \langle \boldsymbol{\alpha}^* - \boldsymbol{\alpha}^{(t)}, \boldsymbol{\Delta}^{(t)} \rangle + \frac{\eta^2}{2} \|\boldsymbol{\Delta}^{(t)}\|_\infty^2
\end{equation}
Summing both sides over $t = 1, \dots, T$:
\begin{equation}
\sum_{t=1}^T \langle \boldsymbol{\alpha}^* - \boldsymbol{\alpha}^{(t)}, \boldsymbol{\Delta}^{(t)} \rangle \le \frac{D_{\Phi}(\boldsymbol{\alpha}^*, \boldsymbol{\alpha}^{(1)})}{\eta} + \frac{\eta T M^2}{2}
\end{equation}
Since $\boldsymbol{\alpha}^{(1)}$ is uniform, $D_{\Phi}(\boldsymbol{\alpha}^*, \boldsymbol{\alpha}^{(1)}) \le \ln D$. Substituting $\eta = \sqrt{\frac{2 \ln D}{M^2 T}}$:
\begin{equation}
R_T \le \frac{\ln D}{\sqrt{\frac{2 \ln D}{M^2 T}}} + \frac{T M^2 \sqrt{2 \ln D}}{2 M \sqrt{T}} = \frac{M \sqrt{T \ln D}}{\sqrt{2}} + \frac{M \sqrt{T \ln D}}{\sqrt{2}} = M \sqrt{2 T \ln D}
\end{equation}
Dividing by $T$ yields an average dual regret of $\mathcal{O}\left( \sqrt{\frac{\ln D}{T}} \right) \to 0$ as $T \to \infty$.
\end{proof}

\begin{theorem}[Immunity to Domain Irreducible Entropy Divergence]
Let domain $k$ possess irreducible noise entropy $H_k \gg 0$, such that $L_k(\theta) = H_k + \delta_k(\theta)$ for any model $\theta$, where $\delta_k(\theta)$ is the reducible model modeling error. Under standard Group DRO optimizing raw loss $L_k$, domain $k$ receives disproportionate weight $\alpha_k \to 1$ as $H_k \to \infty$. Under DoReMi excess loss, the influence of $H_k$ cancels completely:
{\boldmath
\begin{equation}
\Delta_k \equiv \delta_k(\theta_{\text{proxy}}) - \delta_k(\theta_{\text{ref}})
\end{equation}
}
preventing noisy domain takeover.
\end{theorem}

\begin{proof}
By definition of irreducible Bayes entropy $H_k$, the cross-entropy loss decomposes into $L_k(\theta) = \mathbb{E}_{x \sim \mathcal{D}_k}[-\ln p^*(x)] + D_{\text{KL}}(p^*(x) \,\|\, p_\theta(x)) = H_k + \delta_k(\theta)$.
For standard DRO, the objective is $\max_d L_d(\theta) = \max_d (H_d + \delta_d(\theta))$. If $H_k - H_j > \sup_\theta |\delta_j(\theta) - \delta_k(\theta)|$, then domain $k$ dominates the maximum irrespective of model performance.
Under DoReMi:
\begin{equation}
\Delta_k = L_k(\theta_{\text{proxy}}) - L_k(\theta_{\text{ref}}) = \left(H_k + \delta_k(\theta_{\text{proxy}})\right) - \left(H_k + \delta_k(\theta_{\text{ref}})\right) = \delta_k(\theta_{\text{proxy}}) - \delta_k(\theta_{\text{ref}})
\end{equation}
The constant $H_k$ cancels additively, ensuring $\Delta_k$ depends purely on the reducible gap between proxy and reference models.
\end{proof}

\section{Architectural Questions and Real-World Trade-Offs}

\subsection{Architectural Question 1: Reference Model Selection Strategy}
\textbf{Question:} What model should serve as the reference model $\theta_{\text{ref}}$ in DoReMi?\\
\textbf{Answer:} An off-the-shelf pre-trained foundation model of moderate capacity (or an initial checkpoint trained on an equal-weighted mixture). The reference model does not need to be optimal; it only needs to establish a consistent per-domain baseline entropy floor so that excess loss accurately isolates learnable knowledge from noise.

\subsection{Architectural Question 2: Smoothing Epsilon and Domain Starvation}
\textbf{Question:} Why is the smoothing floor $\epsilon > 0$ strictly necessary in production data pipelines?\\
\textbf{Answer:} Without a positive floor ($\epsilon = 0$), if a domain temporarily exhibits zero or negative excess loss during early steps, its weight exponentially decays toward 0. In subsequent training, the proxy model receives zero samples from that domain, making it impossible to observe if the model later requires that domain's data. A floor $\epsilon \approx 0.001$ preserves a non-zero sampling probability, maintaining continuous gradient feedback across all corpora.

\section{Algorithmic Implementation Specification}

\begin{algorithm}[H]
\caption{DoReMi Dynamic Domain Reweighting}
\begin{algorithmic}[1]
\Require Proxy losses $\mathbf{L}_{\text{proxy}} \in \mathbb{R}^D$, reference losses $\mathbf{L}_{\text{ref}} \in \mathbb{R}^D$, current weights $\boldsymbol{\alpha} \in \Delta^D$, step size $\eta$, smoothing $\epsilon$
\Ensure Updated weights $\boldsymbol{\alpha}' \in \Delta^D$, excess losses $\boldsymbol{\Delta} \in \mathbb{R}^D$, worst domain index $d^*$, max excess loss $\Delta^*$
\State Initialize $\boldsymbol{\Delta} \gets [\,]$
\For{$d = 0$ \textbf{to} $D - 1$}
    \State $\Delta_d \gets L_{\text{proxy}, d} - L_{\text{ref}, d}$
    \State Append $\Delta_d$ to $\boldsymbol{\Delta}$
\EndFor
\State $d^* \gets \operatorname{argmax}_d(\boldsymbol{\Delta}), \quad \Delta^* \gets \max(\boldsymbol{\Delta})$
\State Initialize unnormalized vector $\tilde{\boldsymbol{\alpha}} \gets [\,]$
\For{$d = 0$ \textbf{to} $D - 1$}
    \State $\text{clipped\_delta} \gets \min(20.0, \max(-20.0, \eta \cdot \Delta_d))$
    \State Append $\alpha_d \cdot \exp(\text{clipped\_delta})$ to $\tilde{\boldsymbol{\alpha}}$
\EndFor
\State $S \gets \sum_{d=0}^{D-1} \tilde{\alpha}_d$
\State Initialize smoothed vector $\boldsymbol{\alpha}_{\text{smooth}} \gets [\,]$
\For{$d = 0$ \textbf{to} $D - 1$}
    \State Append $(\tilde{\alpha}_d / S) + \epsilon$ to $\boldsymbol{\alpha}_{\text{smooth}}$
\EndFor
\State $S_{\text{final}} \gets \sum_{d=0}^{D-1} \alpha_{\text{smooth}, d}$
\State $\boldsymbol{\alpha}' \gets [ a / S_{\text{final}} \text{ for } a \in \boldsymbol{\alpha}_{\text{smooth}} ]$
\State \Return $(\boldsymbol{\alpha}', \boldsymbol{\Delta}, d^*, \Delta^*, D)$
\end{algorithmic}
\end{algorithm}

\section{Big-$\Theta$ Complexity Analysis}

\begin{itemize}
    \item \textbf{Excess Loss Calculation}: $\Theta(D)$ scalar subtractions.
    \item \textbf{Exponentiated Weight Update}: $\Theta(D)$ exponent and multiplication operations.
    \item \textbf{Simplex Projection and Normalization}: $\Theta(D)$ sum and division passes.
    \item \textbf{Total Time Complexity}: $\Theta(D)$ strictly linear with respect to domain count.
    \item \textbf{Space Complexity}: $\Theta(D)$ auxiliary storage for excess losses and new weights.
\end{itemize}

\section{Single Canonical Repository Reference}

The complete deterministic scalar implementation, verification suite, and unit test benchmarks are published at:
\begin{center}
\url{https://github.com/Chief-Strategist-J/llm-obs-infra}
\end{center}

\end{document}
